% =====================================================================
\BookChapter[R]{R: A Language Built for Statistics}
\label{ch:r}

\begin{researchquestion}
How do you get from a table of measurements to a conclusion that statistics can support?
\end{researchquestion}
\OrientationNote
% =====================================================================

\begin{factsheet}{R}
First appeared & 1993 (version 1.0 in 2000); based on the S language (Bell Labs, 1976)\\
Designed by & Ross Ihaka and Robert Gentleman, University of Auckland; now the R Core Team\\
Paradigms & functional, array (vectorised), object-oriented (S3, S4, R6), statistical\\
Typing & dynamic\\
Execution & interpreted (bytecode compiled)\\
File extensions & \file{.R} (scripts), \file{.Rmd} and \file{.qmd} (R Markdown / Quarto reports), \file{.rds} (data)\\
Cost & free and open source (GPL)\\
Used for & statistics, design of experiments, quality control, reliability analysis,
  data analysis and visualisation, bioinformatics, epidemiology, environmental and
  hydrological modelling, econometrics\\
Famous software & ggplot2, the tidyverse (dplyr, tidyr, readr), Shiny, Bioconductor,
  the CRAN repository of more than 20\,000 packages\\
\end{factsheet}

R is a language and environment for statistical computing and graphics. Where \py is a
general-purpose language that became good at data analysis, R was designed by
statisticians, for statistics, from the start. Almost every statistical method ever
published is available in R, often written by the statisticians who invented it. For
engineers, R is the natural tool for designed experiments, quality and reliability
analysis, environmental data, and any work where careful statistics matters more than raw
speed.

\LearningOutcomes{Create and index vectors in R; handle missing values; store measurements in a data frame and summarise them by group; describe a statistical model with a formula; run a t-test and a linear regression and read their output; and decide when R is a better choice than Python.}

\section{History and philosophy}

\begin{history}
In 1976 John Chambers and colleagues at Bell Labs (the same laboratory that produced C
and Unix) created S, a language for interactive data analysis. S became a commercial
product (S-PLUS). In 1993 two statisticians at the University of Auckland, Ross Ihaka
and Robert Gentleman, wrote a free implementation of a language much like S for teaching,
and named it R --- partly after their first names, partly as a play on S. They released it
under the GNU General Public License in 1995. A worldwide ``R Core Team'' has developed it
since, and the Comprehensive R Archive Network (CRAN) collects contributed packages.
\end{history}

R's design reflects its purpose. The fundamental data structure is the \emph{vector};
missing data are part of the language; statistical models are described by
\emph{formulas} such as \code{y \textasciitilde\ x}; and the \emph{data frame} --- a table
whose columns may have different types --- was R's invention long before pandas copied it.
R is also strongly influenced by the functional language Scheme: functions are values,
and most operations return new objects rather than modifying old ones.

\section{Getting started}

Install R from \url{https://cran.r-project.org} and then \textbf{RStudio} (or
\textbf{Positron}), the IDE that almost everyone uses, from Posit
(\url{https://posit.co}). R also works in VS Code and in Jupyter notebooks.
\codefile{R}{r/hello.R}{hello world}
\begin{lstlisting}[style=shell]
$ Rscript hello.R
\end{lstlisting}
\outfile{r/hello.out}
The \code{[1]} in front of printed output means ``the first element of this vector is
shown here''; long vectors are printed over several lines, each starting with the index of
its first element. \code{cat} prints without quotes and without the index.

\section{The language in brief}

\subsection*{Every value is a vector}

R's traditional assignment operator is the arrow \code{<-} (\code{=} also works). Comments
start with \code{\#}. There are no scalars: a single number is a vector of length one. The
function \code{c()} (``combine'') creates vectors. Here is an interactive session, with
R's output shown after each command:
\begin{lstlisting}[language=R]
> x <- c(276.4, 281.9, 268.7, 273.5, 279.2)   # yield strengths, MPa
> x[1]                      # indices start at 1
[1] 276.4
> x[2:4]                    # 2:4 is the vector 2, 3, 4 (end included)
[1] 281.9 268.7 273.5
> x[-1]                     # a NEGATIVE index REMOVES elements
[1] 281.9 268.7 273.5 279.2
> x[x > 275]                # logical indexing
[1] 276.4 281.9 279.2
> mean(x); sd(x); range(x)
[1] 275.94
[1] 5.116933
[1] 268.7 281.9
> x - mean(x)               # vectorised: the scalar is "recycled"
[1]  0.46  5.96 -7.24 -2.44  3.26
\end{lstlisting}

\begin{pitfall}[Common pitfall: negative indices]
In \py, \code{x[-1]} is the \emph{last} element. In R, \code{x[-1]} is \emph{everything
except the first} element. The last element in R is \code{x[length(x)]}. This is a
frequent source of confusion when moving between the two languages.
\end{pitfall}

\subsection*{Missing measurements are part of the language}

Real measurement data have gaps: a sensor failed, a specimen broke in the grips. R has a
special value \code{NA} (``not available'') for this, and statistical functions treat it
honestly:
\begin{lstlisting}[language=R]
> y <- c(1.2, NA, 3.4)
> mean(y)
[1] NA
> mean(y, na.rm = TRUE)     # explicitly ignore missing values
[1] 2.3
\end{lstlisting}
The mean of a set containing an unknown number is unknown, so R returns \code{NA} unless
you say otherwise. This prevents the silent errors that occur when a missing value is
recorded as zero.

\subsection*{Six data structures cover almost everything}

\begin{center}
\small
\begin{tabularx}{\linewidth}{@{}l>{\raggedright\arraybackslash}X@{}}
\toprule
\textbf{Structure} & \textbf{Description}\\
\midrule
vector & elements of one type: \code{numeric} (double), \code{integer} (\code{5L}),
  \code{character}, \code{logical} (\code{TRUE}/\code{FALSE}), \code{complex}\\
factor & a categorical variable with a fixed set of levels (e.g.\ material, batch, operator)\\
matrix, array & a vector with dimensions; column-major, like Fortran\\
list & an ordered collection of objects of any types (like a \py list or dictionary)\\
data frame & a list of equal-length columns: R's central table structure\\
function & a function is an object like any other\\
\bottomrule
\end{tabularx}
\end{center}
Matrix algebra uses \code{\%*\%} for the matrix product (\code{*} is element-wise),
\code{t()} for transpose and \code{solve(A, b)} for linear systems.

\subsection*{Functions return their last value and work on whole vectors}

R has \code{if}/\code{else}, \code{for}, \code{while} and \code{repeat} with braces, as in
C. Functions are created with \code{function} and assigned to a name; the value of the last
expression is returned automatically.
\begin{lstlisting}[language=R]
reynolds <- function(rho, v, D, mu = 1.0e-3) {
  rho * v * D / mu
}
regime <- function(Re) {
  if (Re < 2300) "laminar" else if (Re < 4000) "transitional" else "turbulent"
}
Re <- reynolds(998, c(0.01, 0.05, 1.0), 0.05)   # vectorised: three Reynolds numbers
sapply(Re, regime)                              # apply a function to each element
# [1] "laminar"      "transitional" "turbulent"
\end{lstlisting}
Explicit \code{for} loops are fine in R, but idiomatic R code prefers vectorised
operations and the \code{apply} family (\code{sapply}, \code{lapply}, \code{vapply},
\code{mapply}), which apply a function to every element of a vector or list.

\section{Case study: the falling skydiver}

\codefile{R}{r/drag.R}{Euler solution of the drag equation}
\outfile{r/drag.out}
\code{sprintf} is vectorised: given vectors of values it produces a vector of formatted
lines, which \code{cat} prints. \code{seq\_len(n)} produces $1, \ldots, n$ and is safer than
\code{1:n}, which gives the unexpected sequence \code{1 0} when $n = 0$.

\section{Data frames and statistics}

R's real power appears when you work with tables of data and statistical models. The
following script reads the tensile test results, summarises them by material, tests
whether the two materials differ, fits a linear model and saves a box plot.
\codefile{R}{r/tensile.R}{data frames, a t-test and a linear model}
\outfile{r/tensile.out}
A few lines of R gave a statistical answer to a typical engineering question (here for
our invented data set). The Welch
t-test shows that the mean yield strengths of the two materials differ significantly
($p < 0.001$). The linear model shows that, within the tested range, the small variations
in specimen diameter explain none of the scatter in the aluminium results ($R^2 \approx 0$):
the scatter is material variability, not a measurement artefact. Notice the
\term{formula} notation \code{yield\_MPa \textasciitilde\ material}: ``yield strength
explained by material''. The same notation is used for t-tests, analysis of variance
(\code{aov}), linear and nonlinear regression (\code{lm}, \code{nls}), generalised linear
models (\code{glm}) and plotting.

\section{The tidyverse and ggplot2}

Many R users work with the \textbf{tidyverse}, a coherent set of packages for data
handling designed by Hadley Wickham and colleagues. Operations are chained with the
\term{pipe} operator \code{|>} (built in since R 4.1; the tidyverse also uses
\code{\%>\%}), which passes the result of one step into the next, so code reads from top to
bottom like a recipe:
\begin{lstlisting}[language=R]
library(tidyverse)
read_csv("tensile.csv") |>
  mutate(area_mm2 = pi * diameter_mm^2 / 4,
         load_kN  = yield_MPa * area_mm2 / 1000) |>
  group_by(material) |>
  summarise(n = n(), mean_load = mean(load_kN), sd_load = sd(load_kN))
\end{lstlisting}
\textbf{ggplot2} implements the ``grammar of graphics'': a plot is built from data, a
mapping of variables to visual properties (\emph{aesthetics}), and layers of geometric
objects. It produces publication-quality figures with very little code:
\begin{lstlisting}[language=R]
ggplot(df, aes(x = material, y = yield_MPa, colour = material)) +
  geom_boxplot(outlier.shape = NA) +
  geom_jitter(width = 0.1) +
  labs(y = "Yield strength (MPa)", x = NULL) +
  theme_bw()
ggsave("yield.pdf", width = 4, height = 3)
\end{lstlisting}
The \py library \textbf{plotnine} copies the ggplot2 syntax, and \textbf{seaborn} offers a
similar statistical-graphics style.

\section{R for engineers}

Some areas where R is particularly strong for engineering work:
\begin{itemize}
\item \textbf{Design of experiments} (packages \code{DoE.base}, \code{FrF2}, \code{rsm}):
  factorial and fractional factorial designs, response surfaces;
\item \textbf{Quality control} (\code{qcc}, \code{SixSigma}): control charts, process
  capability indices;
\item \textbf{Reliability and survival analysis} (\code{survival}, \code{WeibullR}): Weibull
  fits of failure times, life data analysis;
\item \textbf{Time series} (\code{forecast}, \code{fable}): forecasting loads, demand and
  sensor signals;
\item \textbf{Spatial and environmental data} (\code{sf}, \code{terra}): maps, hydrology,
  geostatistics (kriging);
\item \textbf{Reproducible reports}: R Markdown and Quarto documents mix text, code and
  results and render to PDF, Word or HTML --- a test report that regenerates itself when
  new data arrive. (Quarto also works with \py and Julia.)
\item \textbf{Interactive web apps}: \textbf{Shiny} turns an R analysis into a web
  application with sliders and plots, without writing JavaScript.
\end{itemize}

\section{R and Python}

R and \py are often presented as rivals, but many data scientists use both. The
\code{reticulate} package calls \py from R, and \code{rpy2} calls R from \py; Jupyter and
Quarto support both; and data are easily exchanged through CSV, Parquet or Arrow files.
As a rule of thumb: for statistics, designed experiments and statistical graphics, R is
often more convenient; for general programming, machine learning and integration with
other software, \py is usually the better choice.

\section{Strengths and weaknesses}

\begin{center}
\small
\begin{tabularx}{\linewidth}{@{}>{\raggedright\arraybackslash}X>{\raggedright\arraybackslash}X@{}}
\toprule
\textbf{Strengths} & \textbf{Weaknesses}\\
\midrule
Unmatched range of statistical methods & Unusual syntax and many inconsistencies\\
Excellent graphics (ggplot2) & Slow loops; large data must fit in memory\\
Data frames, factors and \code{NA} built in & Less suited to general-purpose programming\\
Reproducible reports (R Markdown, Quarto) and Shiny apps & Several competing styles (base R, tidyverse, data.table)\\
Free; strong academic community & Fewer engineering-specific (non-statistical) packages\\
\bottomrule
\end{tabularx}
\end{center}

\section{Further reading}
\begin{itemize}
\item H.~Wickham, M.~\c{C}etinkaya-Rundel and G.~Grolemund, \emph{R for Data Science},
  2nd ed., O'Reilly, 2023 (free online at \url{https://r4ds.hadley.nz}).
\item N.~Matloff, \emph{The Art of R Programming}, No Starch Press, 2011.
\item D.~C.~Montgomery, \emph{Design and Analysis of Experiments}, Wiley --- the classic
  engineering statistics text, with many R-based supplements.
\end{itemize}

\begin{exercises}
\item In R, compute the mean, median, standard deviation and coefficient of variation of
  the values in \file{yield.txt} (\code{x <- scan("../data/yield.txt")}).
\item Create a vector of 1000 normally distributed random numbers with mean 275 and
  standard deviation 5 (\code{rnorm}). Plot a histogram and add the theoretical density
  curve.
\item What do \code{c(1, 2, 3) + c(10, 20, 30, 40, 50, 60)} and
  \code{c(1, 2) + c(10, 20, 30)} return? Explain R's recycling rule. Why is the second
  case accompanied by a warning?
\item Use \code{lm} to fit a straight line to the data $x = (0, 1, 2, 3, 4)$,
  $y = (0.1, 1.9, 4.2, 5.8, 8.1)$. Report the slope, intercept and $R^2$, and plot the data
  with the fitted line.
\item Rewrite the tensile analysis with the tidyverse and ggplot2.
\end{exercises}
